\section{Discussion}
The central finding is that ADC, despite collapsing the multi-component signal to a single scalar, ranks ROIs for tumor detection almost identically to the optimal eight-component classifier trained on the full spectrum (Figure~\ref{fig:adc_discriminant}).
Rather than an impoverished summary of a richer signal, ADC captures the dominant axis of spectral variation relevant to detection---the spectral decomposition does not obsolete ADC, it explains it.
This is consistent with recent works in which advanced diffusion models matched ADC for prostate cancer detection \citep{langkilde2018evaluation, he2025improved}.

\subsection{Why ADC works}
Both readouts reduce the same eight-bin spectrum to one number, with nearly the same weights on the same two bins (Figure~\ref{fig:sensitivity}).
The discriminant is linear in the fractions by construction. ADC is not, but over the range of spectra this cohort spans it is well approximated by its linearization about the mean spectrum, the sensitivity vector $\partial\text{ADC}/\partial\Rvec$, so both act as linear functionals of the spectrum.
Tumor-versus-normal variation is concentrated in two well-identified compartments, restricted diffusion rising as free water falls, so the spectra vary along an essentially one-dimensional axis and two linear readouts projecting onto it must rank ROIs almost identically.

This resolves an apparent mismatch. The two weight vectors are only moderately anti-correlated bin by bin ($r = -0.79$ and $-0.88$), yet the scores they produce agree at $|r| > 0.97$. They disagree mainly on the intermediate bins, where ADC retains appreciable sensitivity but the discriminant places almost none. What decides the agreement between scores is not how similar the weights are, but how similar they are along the directions in which the spectra actually vary, and those intermediate directions carry too little tumor--normal variation for the disagreement to register.

Discriminative power is sensitivity \emph{times} contrast. The intermediate fractions carry little tumor--normal contrast beyond the restricted and free-water pair and are poorly identified (Figure~\ref{fig:fisher}), so ADC's sensitivity to them is spent on a low-contrast, noisy direction that moves neither ranking. They are redundant rather than uninformative---each co-varies with the dominant axis and discriminates on its own, yet adds nothing once those two compartments are present (Table~\ref{tab:auc}).

This is why resolving more diffusion compartments \citep{panagiotaki2014microstructural, conlin2021improved} need not improve detection over ADC in an acquisition such as ours. The compartments such models add are not independent of those already resolved, and the contrast they carry is largely the same contrast re-expressed, duplicating the restricted-to-free-water axis rather than supplying a second one, which is why our eight-fraction model does not beat the two-fraction one. Where richer models do gain, the gain traces to contrast this acquisition lacks. Hybrid multidimensional MRI separates tumor from normal almost perfectly, but by adding an echo-time dimension, not finer diffusion compartments \citep{chatterjee2018hybrid}.
One-dimensional observable variation is not in tension with the ill-conditioning of the design. The singular values of $\Umat$ span five orders of magnitude, so at the cohort-median noise level only five of the eight orthogonal directions in spectrum space are estimated with a standard deviation below $0.1$ in fraction units, the remaining three exceeding $0.6$. Over $90\%$ of the tumor--normal difference vector nevertheless lies inside that resolvable subspace. Ill-conditioning caps how many spectral directions the acquisition can resolve, but what caps detection is the single contrast direction biology supplies.

\subsection{Grade-dependent spectral shifts}
% NOTE (2026-07-26): kept descriptive. The spectrum makes grade-associated biology visible, but we claim no per-ROI grading advantage over ADC and report no grading AUC (small n). The formal discriminative-power framework was removed per Stephan's review; it is deferred to a follow-up paper.
Any signal beyond ADC would have to live in the intermediate bins, where ADC is least sensitive, and this is plausibly where the decomposition adds value (Figure~\ref{fig:spectrum_ggg}).
The free-water bin ($D = 3.0\;\umsmm$, a proxy for intact gland lumen) collapses as soon as tumor is present ($0.39$ normal to $0.21$ at GGG\,$=\,1$) and then changes little ($0.16$ at GGG\,$\geq 3$), separating tumor from normal rather than grade from grade.
By contrast the restricted-cellular bin ($D = 0.25\;\umsmm$) rises monotonically with grade ($0.07$ normal to $0.31$ at GGG\,$\geq 3$), and the glandular-epithelial bin ($D = 2.0\;\umsmm$) is preserved in low-grade tumor and depleted at higher grades ($0.26$ normal, $0.28$ at GGG\,$=\,1$, $0.17$ at GGG\,$\geq 3$).

Because the grade groups are too small to subdivide, this comparison pools the two zones, whose normal baselines differ (Results), and the graded tumors are predominantly peripheral-zone (21 of 29) while the normal baseline is evenly split. Two checks indicate that pooling does not manufacture the pattern. First, it is conservative. Peripheral-zone normal tissue carries more free water ($0.48$ vs.\ $0.30$) and less restricted mass ($0.06$ vs.\ $0.09$) than transition-zone tissue, so a zone-matched baseline would show a \emph{larger} free-water collapse and a \emph{larger} restricted rise. Second, the pattern survives restriction to the peripheral zone alone ($n = 21$), where the restricted bin still rises across GGG\,$=\,1$, $2$ and $\geq 3$ ($0.16 \to 0.17 \to 0.28$, against $0.06$ in normal tissue) and the glandular-epithelial bin still falls ($0.29 \to 0.25 \to 0.17$), its within-zone rank correlation with grade remaining nominally significant (Spearman $\rho = -0.44$, $P = 0.045$).

This follows the cellular reorganization of prostate adenocarcinoma---loss of patent lumina, glandular effacement, rising cell-packing density---and is consistent with histological work relating epithelium, stroma, and lumen volumes to Gleason pattern and to ADC \citep{chatterjee2015changes}, volumes that multidimensional MRI recovers in vivo \citep{chatterjee2022validation}, and with reports that advanced diffusion models improve grade differentiation where ADC plateaus \citep{johnston2019verdict, he2025improved}.
We present this as a candidate interpretation the decomposition makes observable, not a validated finding. The grade groups are small ($n = 8$--$9$) and unbalanced by zone, so no grading AUC is claimed.

\subsection{What the Bayesian posterior contributes}
At tuned regularization MAP and the NUTS posterior mean agreed closely on the well-identified components and gave nearly interchangeable discriminant scores ($r = 0.99$ PZ, $0.97$ TZ), so closed-form MAP is sufficient for classification.
Their divergence on the intermediate components reflects the ill-posedness, not a deficiency of either method. Those fractions are genuinely underdetermined by fifteen b-values, even at the low noise of our endorectal-coil acquisition, and the sampler reproduces the same smearing on simulated ground truth (Figure~\ref{fig:validation}).
The posterior makes the identifiability gradient explicit, and the two-step tightening of the attainable estimation error in Figure~\ref{fig:fisher}b is what moves this classically ill-conditioned separation \citep{jalnefjord2019optimization} into a regime where the well-constrained components can be believed. The posterior's value is therefore upstream of the tumor-versus-normal decision rather than within it.

Consistent with this, propagating the full posterior through the classifier added no diagnostic value beyond the point estimate (Figure~\ref{fig:uncertainty}). Detection rides on the two well-identified compartments, leaving little uncertainty in the discriminant direction to propagate, and supplying the posterior standard deviations as extra features left AUC unchanged.
The credible intervals on $P(\mathrm{tumor})$ are on average $2.4\times$ wider for misclassified ROIs, but this is a geometric property of the logistic link---its slope is largest at the decision boundary---not a contribution of the spectral posterior, and in logit space the difference is no longer significant.
Because errors concentrate near the boundary, distance from it does flag the least reliable predictions, but that is a standard property of the logistic model, available from the point estimate alone and shared by both classifiers.

\subsection{Pixel-wise feasibility}
The framework extends to the voxel level (Figure~\ref{fig:pixelwise}).
On a single peripheral-zone slice the restricted-diffusion map reproduced the ADC pattern, the free-water map was its negative, and the per-voxel discriminant anti-correlated with ADC ($r = -0.96$). The uncertainty map showed the eight-bin detector $1.8\times$ more uncertain than the two-bin one, because it weights the unidentifiable intermediate compartments.
It also exposes a calibration caveat. A single voxel has far lower SNR than an ROI average, so its high-b measurements sit near the noise floor, which mimics slow decay and inflates the restricted fraction. The ROI-trained detector reads that inflation as evidence of tumor and shifts every voxel score toward the tumor class, even though the within-slice spatial contrast is preserved. This is why all quantitative results are reported at the ROI level and these maps are exploratory. The uncertainty does not yet improve detection, but it indicates which compartments are locally reliable.

\subsection{Limitations and future directions}
First, the tumor ROIs were placed on multiparametric MRI, including DWI and ADC maps, introducing a potential circularity. Lesions conspicuous on ADC may be overrepresented \citep{langkilde2018evaluation}, inflating ADC's apparent performance relative to spectral features. Compounding this, the detection-positive ROIs are radiological targets rather than biopsy-confirmed cancers---11 of the 40 were not confirmed malignant---so the detection contrast is partly against the radiologist's own ADC-informed reading. Only the Gleason grade analysis, from biopsy pathology, is free of this circularity, though its graded sample ($n = 29$) is too small for a reliable between-method comparison.
Second, the equivalence is specific to ROI-level detection of delineated lesions. For patient-level, delineation-free detection, restriction-spectrum imaging outperforms whole-gland minimum ADC \citep{domingo2025restriction}, a different task.
Third, the diffusivity grid is fixed. A finer or free-form grid might reveal more structure at the cost of greater ill-posedness.
Fourth, a single echo-time and diffusion-time configuration was analyzed. Shorter timings enabled by recent gradient hardware \citep{zhu2024human, elsaid2026nonlinear} could change both the SNR and the relative weight of the visible compartments.
Finally, the pixel-wise result rests on a single patient and requires larger studies with histopathological correlation.

The framework applies to other organs in which multi-b DWI is acquired, and the Fisher-information analysis could guide joint optimization of b-value protocols and diffusivity grids. Integration with quantitative T\textsubscript{2} mapping, which independently resolves luminal water \citep{sabouri2017luminal}, could separate the grading axis from the detection axis characterized here.
